%=======================================================================
\chapter{Transport Formulas}
%=======================================================================

The balance laws of the next chapter are statements about \emph{integrals}:
how much mass, momentum or energy a lump of material carries, and how fast
that amount changes. Every one of them therefore needs an answer to a
single question:
\[
   \frac{\dl}{\dl t}\int_{P_{t}}(\,\cdot\,)\,\dl v=\ ?
\]
The question is harder than it looks because \emph{both} the integrand and
the region of integration depend on time. This chapter answers it, in four
settings of increasing generality: a material region (Reynolds' formula),
a moving surface (the flux derivative), a region that is not material at
all, and a region across which the integrand jumps. Nothing here is
physics yet; it is the calculus the physics will be written in.

Chapter 1 already proved the divergence theorem and the two variants that
follow from it, Theorem~\ref{thm:div} and Corollary~\ref{cor:divforms}. We
do not prove them again. Section~\ref{sec:divvar} instead collects them in
the form the balance laws will use and then makes a point Chapter 1 had no
reason to make: two of the three are \emph{symbolic} statements, and it
matters to know exactly what they mean.

%=======================================================================
\section{The divergence theorem and its variants}\label{sec:divvar}
%=======================================================================

For a vector field $\ba(\bx)$ on a volume $R\subset\Et$ bounded by a
piecewise smooth surface $\partial R$, with $\bn$ the piecewise continuous
exterior unit normal, Theorem~\ref{thm:div} is
\begin{equation}
   \int_{R}\dvg\ba\,\dl v=\int_{\partial R}\ip\ba\bn\,\dl a ,
   \tag{5.1}\label{eq:5.1}
\end{equation}
or, in Cartesian coordinates and components,
\begin{equation}
   \int_{R}a_{i,i}\,\dl v=\int_{\partial R}a_{i}n_{i}\,\dl a .
   \tag{5.2}\label{eq:5.2}
\end{equation}

Both of the variants we need come from \eqr{5.1} by the same device: feed
it a field built from an \emph{arbitrary constant vector} $\bc$, and then
let the arbitrariness of $\bc$ do the work. It is worth isolating the
device, because it is used four times in this chapter and a dozen more in
the next.

\begin{lemma}[the constant-vector argument]\label{lem:cvec}
Let $\bu$ be a fixed vector and suppose $\ip\bc\bu=0$ for every vector
$\bc$. Then $\bu=\bze$.
\end{lemma}

\begin{proof}
The hypothesis holds for \emph{every} $\bc$, so in particular for
$\bc=\bu$, giving $\ip\bu\bu=|\bu|^{2}=0$; positive definiteness of the
inner product, Definition~\ref{def:ip}(d), then forces $\bu=\bze$.
\end{proof}

\begin{remark}[why one may not simply ``cancel $\bc$'']\label{rem:cancel}
The step is worth stating as a lemma because the tempting shortcut is
wrong. From $\ip\bc\bu=\ip\bc\bw$ one may \emph{not} conclude $\bu=\bw$
for a \emph{particular} $\bc$: any $\bu-\bw$ perpendicular to that $\bc$
satisfies it. It is the quantifier ``for all $\bc$'' that carries the
argument, and the choice $\bc=\bu$ is what converts it into a statement
about a norm. The same pattern will reappear in the localization theorem,
where ``for all subvolumes'' plays the role that ``for all $\bc$'' plays
here.
\end{remark}

\subsection*{The gradient of a scalar}

Let $\varphi(\bx)$ be a smooth scalar field and put $\ba=\varphi\bc$ with
$\bc$ spatially uniform. The choice looks arbitrary; it is not. We want a
statement about $\grad\varphi$, and \eqr{5.1} only knows how to handle a
divergence, so we need a vector field whose divergence produces
$\grad\varphi$. Multiplying $\varphi$ by a constant vector is the simplest
thing that does it, and keeping the vector constant is what makes the
product rule collapse to a single term.

In components $a_{i}=\varphi c_{i}$. Differentiate, remembering that the
$c_{i}$ are numbers and not functions of position, so they pass through
$\partial/\partial x_{i}$ untouched:
\[
   \dvg\ba=a_{i,i}=(\varphi c_{i})_{,i}=\varphi_{,i}c_{i}+\varphi\,c_{i,i}
   =\varphi_{,i}c_{i},
\]
the last term vanishing because $c_{i,i}=0$ for a uniform field. On the
boundary,
\[
   \ip\ba\bn=a_{i}n_{i}=\varphi c_{i}n_{i}.
\]
Put both into \eqr{5.2}:
\[
   \int_{R}\varphi_{,i}c_{i}\,\dl v=\int_{\partial R}\varphi c_{i}n_{i}\,\dl a .
\]
The $c_{i}$ are constants, so they may be taken outside the integral
signs. Doing that on both sides and moving everything to one side,
\[
   c_{i}\left[\int_{R}\varphi_{,i}\,\dl v
              -\int_{\partial R}\varphi n_{i}\,\dl a\right]=0 .
\]
The bracket depends on $i$ but not on $\bc$, so the display says
$\ip\bc\bu=0$ where $\bu=u_{i}\be_{i}$ has components
\begin{equation}
   u_{i}=\int_{R}\varphi_{,i}\,\dl v-\int_{\partial R}\varphi n_{i}\,\dl a .
   \tag{5.3}\label{eq:5.3}
\end{equation}
As $\bc$ is arbitrary, Lemma~\ref{lem:cvec} gives $\bu=\bze$, i.e.
\begin{equation}
   \int_{R}\varphi_{,i}\,\dl v=\int_{\partial R}\varphi n_{i}\,\dl a ,
   \tag{5.4}\label{eq:5.4}
\end{equation}
valid for every smooth scalar field $\varphi$.

\subsection*{The divergence of a tensor}

Let $\bA(\bx)$ be a smooth tensor field and put $\ba=\bA^{t}\bc$, again
with $\bc$ uniform. Then $a_{j}=(\bA^{t}\bc)_{j}=A_{ij}c_{i}$, so
\[
   \dvg\ba=a_{j,j}=c_{i}A_{ij,j},
   \qquad
   \ip\ba\bn=a_{j}n_{j}=c_{i}A_{ij}n_{j},
\]
and \eqr{5.2} yields $\ip\bc\bu=0$ with
\begin{equation}
   u_{i}=\int_{R}A_{ij,j}\,\dl v-\int_{\partial R}A_{ij}n_{j}\,\dl a .
   \tag{5.5}\label{eq:5.5}
\end{equation}
The arbitrariness of $\bc$ again gives $\bu=\bze$, hence
\begin{equation}
   \int_{R}A_{ij,j}\,\dl v=\int_{\partial R}A_{ij}n_{j}\,\dl a
   \tag{5.6}\label{eq:5.6}
\end{equation}
for every smooth tensor field. We may write \eqr{5.4} and \eqr{5.6} in the
symbolic forms
\begin{equation}
   \int_{R}\grad\varphi\,\dl v=\int_{\partial R}\varphi\,\bn\,\dl a
   \qquad\text{and}\qquad
   \int_{R}\dvg\bA\,\dl v=\int_{\partial R}\bA\bn\,\dl a ,
   \tag{5.7}\label{eq:5.7}
\end{equation}
which are Corollary~\ref{cor:divforms}. There is an important caveat
attached to them, and it is the subject of the next few paragraphs.

\subsection*{What ``symbolic'' means here}

\begin{remark}[the integrals in \eqr{5.7} are not integrals of
vectors]\label{rem:symbolic}
An integral is a limit of sums, and to add vectors that live at different
points of $R$ one must first agree on how to compare them. In $\Et$ with a
\emph{fixed} orthonormal basis there is such an agreement, because parallel transport is trivial there,
and that is exactly what \eqr{5.3} and \eqr{5.5}
use when they integrate \emph{components}. The symbolic forms \eqr{5.7}
are shorthand for those component statements. They are not independent
definitions, and they are not valid for an arbitrary basis field.
\end{remark}

The point deserves the check the book gives it. What we must verify is
that the object $\bu$ built in \eqr{5.3} really is a vector: that its
components in two fixed orthonormal bases are related the way components
of a vector are. Let $\{\be_{i}'\}$ be a second fixed orthonormal basis
and $a_{ij}=\ip{\be_{i}}{\be_{j}'}$, so that by Section~\ref{ss:cob} the
rule to be confirmed is $u_{i}=a_{ij}u_{j}'$. With
$\bx=x_{i}\be_{i}=x_{j}'\be_{j}'$, hence $x_{j}'=a_{ij}x_{i}$, and with
$\varphi_{,j}'=\partial\varphi/\partial x_{j}'$, the chain rule turns
\eqr{5.3} into
\begin{equation}
   u_{i}=\int_{R}a_{ij}\varphi_{,j}'\,\dl v
        -\int_{\partial R}a_{ij}n_{j}'\,\dl a
        =a_{ij}\Big(\int_{R}\varphi_{,j}'\,\dl v
        -\int_{\partial R}\varphi n_{j}'\,\dl a\Big)
        =a_{ij}u_{j}' ,
   \tag{5.8}\label{eq:5.8}
\end{equation}
as required. The whole step turns on being able to move $a_{ij}$ across
the integral sign, and that is legitimate precisely because the $a_{ij}$
are \emph{constants}. Consequently the statement $\bu=\bze$ is basis
independent: $u_{i}'=0$ if and only if $u_{i}=0$.

\begin{remark}[where it fails]\label{rem:polarfail}
Replace the fixed basis by the polar frame of \eqr{1.122}. Then the
matrix relating the two sets of components depends on position, cannot be
taken outside the integral, and \eqr{5.8} collapses. So an equation like
$\int\grad\varphi\,\dl v=\int\varphi\bn\,\dl a$ read with $\grad\varphi$
and $\bn$ resolved on $\{\ber,\bet,\be_{3}\}$ is simply false. Equation
\eqr{5.1}, by contrast, needs no caveat at all: its integrands
$\dvg\ba$ and $\ip\ba\bn$ are scalars, and a scalar is the same number in
every basis.
\end{remark}

%-----------------------------------------------------------------------
\subsection{Example: The Piola identity}\label{sec:piolaid}
%-----------------------------------------------------------------------

Here is the first use, and a good one: a purely kinematical identity
falls out of \eqr{5.7} with the most trivial integrand imaginable.

\begin{remark}[why one feeds an identity its most boring
input]\label{rem:whyphione}
The choice $\varphi=1$ below looks like it must give nothing, and that
impression is worth correcting, because the move is standard and will be
used again.

A general identity holds for \emph{every} admissible input. Its content is
therefore a whole family of statements, and one is free to pick the member
that says the most. Picking a complicated $\varphi$ makes both sides
complicated and relates two hard things; picking $\varphi=1$ annihilates the left side, since $\grad 1=\bze$, and
leaves the right side standing alone with nothing to be balanced
against. A statement of the form ``this expression
equals zero'' is far stronger than ``this expression equals that one'', so
the most trivial input yields the sharpest conclusion. This is the same
logic by which $\bc=\bu$ was the useful choice in Lemma~\ref{lem:cvec},
and by which $\varphi=1$ will give \eqr{5.16} in the next section.

The general principle: to extract something specific from something
general, look for the input that makes one side collapse.
\end{remark}

Let $S\subset B$ be an arbitrary part of the body, occupying $P_{t}\subset
\kappa_{t}$ now and $P\subset\kappa$ in the reference configuration, so
that $P_{t}=\bchi_{\kappa}(P,t)$. Apply \eqr{5.7}$_{1}$ on $P_{t}$ with
$\varphi=1$. Its gradient vanishes, so the left side is zero and
\begin{equation}
   \bze=\int_{\partial P_{t}}\bn\,\dl a
       =\int_{\partial P}\bF^{*}\bN\,\dl A
       =\int_{P}\Dvg\bF^{*}\,\dl V ,
   \tag{5.9}\label{eq:5.9}
\end{equation}
where the middle step is the Piola--Nanson formula \eqr{2.54},
$\bn\,\dl a=\bF^{*}\bN\,\dl A$, with $\bN$ the exterior unit normal on
$\partial P$ and $\bF^{*}$ the cofactor of the deformation gradient
(Definition~\ref{def:cof}), and the last step is the referential form of
\eqr{5.7}$_{2}$. Here $\Dvg$ differentiates with respect to
$\bX\in\kappa$; in components
\begin{equation}
   \Dvg\bF^{*}=F^{*}_{iA,A}\be_{i}.
   \tag{5.10}\label{eq:5.10}
\end{equation}

What \eqr{5.9} has established is that
\[
   \int_{P}\Dvg\bF^{*}\,\dl V=\bze
\]
for one particular $P$. That alone says nothing pointwise: an integral can
vanish while its integrand does not, as $\int_{-1}^{1}x\,\dl x=0$ shows.
The extra ingredient is that $S$ was an arbitrary part of the body, so $P$
is an arbitrary subvolume of $\kappa$, and the integral vanishes for
\emph{every} one of them at once.

That is exactly the hypothesis of the localization theorem,
Theorem~\ref{thm:localization}, whose other hypothesis, continuity of the
integrand, holds because $\bF^{*}$ is built polynomially from $\bF$ and
$\bF$ is continuously differentiable. The theorem's argument is worth
recalling in one line, since it is the same shape as
Lemma~\ref{lem:cvec}: if the integrand were non-zero at some point, it
would be non-zero with one sign throughout a small ball around that point
by continuity, and integrating over that ball alone would give a non-zero
answer, contradicting the hypothesis. Freedom to choose the region plays
here the part that freedom to choose $\bc$ played there.

Applying it componentwise gives the \emph{Piola identity}
\begin{equation}
   \boxed{\ \Dvg\bF^{*}=\bze\quad\text{at every }\bX\in\kappa.\ }
   \tag{5.11}\label{eq:5.11}
\end{equation}

\begin{remark}[what has just happened]\label{rem:piolawhat}
Read \eqr{5.9} from left to right and the identity looks like a fact about
integrals; read it from right to left and it is a fact about the
deformation alone, with no integral in sight. That is the pattern of the
whole subject: an integral statement, valid for every part, plus
localization, equals a pointwise statement. Chapter 6 will run the same
argument on mass, momentum and energy. Here the ``physics'' is the empty
statement that a closed surface has zero total normal. Read geometrically,
$\int_{\partial P_{t}}\bn\,\dl a=\bze$ says that a uniform pressure
acting all round a closed surface cannot push the body anywhere. Even
that empty-looking statement yields something: the nine derivatives of $\bF^{*}$ are not independent.
\end{remark}

%=======================================================================
\section{Reynolds' transport formula}\label{sec:reynolds}
%=======================================================================

Let $\varphi$ be a scalar field on the body, the volumetric density of
some quantity, so that $\varphi=\hat\varphi(\bX,t)=\tilde\varphi(\bx,t)$
with $\bx=\bchi_{\kappa}(\bX,t)$. We want
\begin{equation}
   \frac{\dl}{\dl t}\int_{P_{t}}\tilde\varphi(\bx,t)\,\dl v
   \qquad\text{at fixed }S,
   \tag{5.12}\label{eq:5.12}
\end{equation}
the derivative following a fixed set of \emph{material points}. Both the
integrand and the domain move, and there is no rule for differentiating
under an integral sign whose limits are unknown functions of $t$. The
remedy is the one used throughout Chapter 2: pull back to the reference
configuration, where the domain does not move at all.

With $\dl v=J\,\dl V$ and $P$ fixed in time,
\begin{equation}
\begin{aligned}
   \frac{\dl}{\dl t}\int_{P_{t}}\tilde\varphi(\bx,t)\,\dl v
   &=\frac{\dl}{\dl t}\int_{P}\hat\varphi(\bX,t)\,J\,\dl V\\
   &=\int_{P}(\varphi J)^{\textstyle\cdot}\,\dl V\\
   &=\int_{P}\big(\dot\varphi J+\varphi\dot J\big)\,\dl V ,
\end{aligned}
   \tag{5.13}\label{eq:5.13}
\end{equation}
the interchange of $\dl/\dl t$ with $\int_{P}$ being permissible because
$P$ no longer depends on $t$, provided $\varphi J$ and $(\varphi J)^{
\textstyle\cdot}$ are continuous in $\bX$ and $t$ jointly and dominated on
$P$ by an integrable function of $\bX$.

\begin{remark}[the one move this chapter keeps making]\label{rem:pullback}
Step \eqr{5.13}$_{1}$ is the whole method, and it will be made three more
times, so it is worth naming: \emph{when a domain moves, change variables
until it does not, then differentiate}.

The obstacle was never the integrand, since differentiating under an
integral sign is routine. The obstacle was the domain. There is no general rule for
$\frac{\dl}{\dl t}\int_{D(t)}$, and the substitution $\dl v=J\,\dl V$
removes the difficulty entirely by trading a moving region $P_{t}$ for the
fixed region $P$ and paying with an extra factor $J$ inside. All the time
dependence is then in the integrand, where the ordinary rules apply. The
price is that we must know how $J$ changes, and that is exactly what
Problem 2 of Chapter 4 supplied.

This is the same manoeuvre as pulling back to the reference configuration
in Chapter 2, and it is the reason the reference configuration is worth
carrying around at all: not because it is physically distinguished, for it is not and
Remark~\ref{rem:whyL} of Chapter 4 insisted on the point, but because it
holds still. Watch for it again at \eqr{5.20}, where
Piola--Nanson plays the role $J$ plays here and converts a moving surface
into a fixed one.

And watch for where it \emph{fails}: in Section~\ref{sec:nonmat} the
region is not material, so there is no fixed $P$ to pull back to, the
method has nothing to bite on, and we are thrown back on the definition of
the derivative. That is the reason \eqr{5.23}, alone in this chapter, is
proved from first principles.
\end{remark}

Now dispose of the $\dot J$ that appeared in \eqr{5.13}. By Problem 2 of
Chapter 4, $\dot J=J\tr\bL$ with $\bL=\grad\bv$, and
$\tr\bL=\tr(\grad\bv)=\dvg\bv$ by Proposition~\ref{prop:divv}. Substitute
this into the integrand of \eqr{5.13} and factor out the $J$ that both
terms then contain:
\[
   \dot\varphi J+\varphi\dot J
   =\dot\varphi J+\varphi J\dvg\bv
   =J\big(\dot\varphi+\varphi\dvg\bv\big).
\]
This factoring is the point of the whole calculation. The integral now
reads
\[
   \int_{P}\big(\dot\varphi+\varphi\dvg\bv\big)J\,\dl V ,
\]
and $J\,\dl V$ is $\dl v$, which is the substitution of \eqr{5.13}$_{1}$
run backwards. Undoing it carries the integral back to the moving region
and leaves
\begin{equation}
   \frac{\dl}{\dl t}\int_{P_{t}}\varphi\,\dl v
   =\int_{P_{t}}\big(\dot\varphi+\varphi\dvg\bv\big)\,\dl v .
   \tag{5.14}\label{eq:5.14}
\end{equation}
Notice what the round trip accomplished. We went to the reference
configuration only to be allowed to differentiate, and came straight back.
Nothing referential survives in \eqr{5.14}: both sides are integrals over
$P_{t}$ and every symbol in it is spatial. The reference configuration was
scaffolding, put up to make one step legal and taken down immediately
afterwards. That is why the formula holds even for materials for which no
natural reference configuration exists.

This is already a transport formula, and a useful one, but it is written
in terms of the material derivative $\dot\varphi$. For a spatial
description we want $\varphi'$ instead. Substitute
$\dot\varphi=\varphi'+\ip{\grad\varphi}\bv$ from \eqr{2.20} into the
integrand of \eqr{5.14}:
\[
   \dot\varphi+\varphi\dvg\bv
   =\varphi'+\ip{\grad\varphi}\bv+\varphi\dvg\bv .
\]
The last two terms are exactly the two halves of a product rule. By
Problem 24(c) of Chapter 1,
\[
   \dvg(\varphi\bv)=\ip{\grad\varphi}\bv+\varphi\dvg\bv ,
\]
so those two terms collapse into a single divergence and the integrand
becomes $\varphi'+\dvg(\varphi\bv)$. Splitting the integral in two,
\[
   \frac{\dl}{\dl t}\int_{P_{t}}\varphi\,\dl v
   =\int_{P_{t}}\varphi'\,\dl v+\int_{P_{t}}\dvg(\varphi\bv)\,\dl v .
\]
The second integral is a divergence over a volume, which is precisely what
\eqr{5.1} converts into a surface integral, with $\ba=\varphi\bv$:
\[
   \int_{P_{t}}\dvg(\varphi\bv)\,\dl v
   =\int_{\partial P_{t}}\ip{\varphi\bv}\bn\,\dl a
   =\int_{\partial P_{t}}\varphi\,\ip\bv\bn\,\dl a ,
\]
the scalar $\varphi$ coming out of the inner product. Therefore
\begin{equation}
   \boxed{\ \frac{\dl}{\dl t}\int_{P_{t}}\varphi\,\dl v
   =\int_{P_{t}}\varphi'\,\dl v
   +\int_{\partial P_{t}}\varphi\,\ip\bv\bn\,\dl a\ }
   \tag{5.15}\label{eq:5.15}
\end{equation}
This is \emph{Reynolds' transport formula}.

It is worth seeing why the product rule was the right thing to look for.
We had a volume integral and wanted a surface integral, and \eqr{5.1} is
the only tool that makes that conversion. But \eqr{5.1} accepts a
divergence and nothing else, so the whole step amounted to asking whether
the leftover terms could be assembled into one. They could, and not by
luck: $\ip{\grad\varphi}\bv+\varphi\dvg\bv$ is the full expansion of
$\dvg(\varphi\bv)$, so the two terms that appeared were never anything
but a divergence written out. Recognising a product rule run backwards is
the recurring skill in this subject, and it is worth practising on the
simpler scalar case $\frac{\dl}{\dl x}(fg)=f'g+fg'$ until the pattern is
automatic.

\begin{remark}[the two clocks: $\dot\varphi$ against
$\varphi'$]\label{rem:twoclocks}
Three different time derivatives appear between \eqr{5.12} and
\eqr{5.15}, and the whole formula is really a statement about how they
relate. It is worth setting them side by side, since every term of
\eqr{5.15} is an answer to ``which one?''.

\[
\begin{aligned}
&\frac{\dl}{\dl t}\int_{P_{t}}\varphi\,\dl v
  &&\text{follows a fixed lump of \emph{material}; the domain moves,}\\
&\dot\varphi
  &&\text{follows a fixed material \emph{point}; hold }\bX\text{ fixed,}\\
&\varphi'
  &&\text{follows a fixed \emph{place}; hold }\bx\text{ fixed.}
\end{aligned}
\]

The first is what a balance law asks about, since it weighs the same piece
of matter twice. The second is what a particle experiences. The third is what
a fixed probe records. Chapter 2 already tied the last two together in
\eqr{2.20}, $\dot\varphi=\varphi'+\ip{\grad\varphi}\bv$: the difference is
the convective term, the change you meet merely by being carried to a new
place where $\varphi$ happens to differ.

Equation \eqr{5.14} is the transport formula stated in the second clock,
\eqr{5.15} the same content in the third. Neither is more correct; they
differ in what has to be known to use them. If the data come as $\hat\varphi(\bX,t)$, as a constitutive statement
about the material would, then the overdot is the native derivative and
\eqr{5.14} is the shorter road. If the data come
as $\tilde\varphi(\bx,t)$, read off a field in space, the prime is native
and \eqr{5.15} is. The convective term is the toll paid at the border, and
Remark~\ref{rem:xprime} of Chapter 2 is the warning about forgetting to
pay it.
\end{remark}

\begin{remark}[how to read \eqr{5.15}]\label{rem:readreynolds}
The amount of $\varphi$-stuff inside the moving lump changes for exactly
two reasons, and the formula separates them. The volume integral is the
change you would measure standing still at each point of space: it is
local creation. The surface integral is the change that comes from the
boundary sweeping over new material: $\ip\bv\bn$ is the outward normal
speed of the boundary, $\varphi\ip\bv\bn\,\dl a$ the rate at which
$\varphi$-stuff is swept past $\dl a$. Nothing is created there; the lump
merely takes in territory. Equation \eqr{5.14} says the same thing in the
other split, following the material rather than the place. Both are exact;
which one is convenient depends on whether the data are given as
$\hat\varphi$ or $\tilde\varphi$.
\end{remark}

Taking $\varphi=1$, so that $\varphi'=0$, gives the rate of change of a
material volume as the flux of material velocity through its boundary:
\begin{equation}
   \frac{\dl}{\dl t}\int_{P_{t}}\dl v=\int_{\partial P_{t}}\ip\bv\bn\,\dl a .
   \tag{5.16}\label{eq:5.16}
\end{equation}
Comparing with \eqr{5.14} at $\varphi=1$ recovers $\dot J/J=\dvg\bv$ in
integrated form, which is a useful consistency check and also the reason
$\dvg\bv=0$ is the pointwise statement of isochoric motion.

%=======================================================================
\section{Extensions and generalizations}\label{sec:extensions}
%=======================================================================

%-----------------------------------------------------------------------
\subsection{Reynolds' formula for vector fields}\label{sec:reyvec}
%-----------------------------------------------------------------------

Apply \eqr{5.15} to $\varphi=u_{i}$, one Cartesian component of a vector
field $\bu$ resolved on the fixed basis $\{\be_{i}\}$:
\begin{equation}
   \frac{\dl}{\dl t}\int_{P_{t}}u_{i}\,\dl v
   =\int_{P_{t}}u_{i}'\,\dl v
   +\int_{\partial P_{t}}u_{i}\,\ip\bv\bn\,\dl a ;
   \qquad i\in\{1,2,3\}.
   \tag{5.17}\label{eq:5.17}
\end{equation}
Multiply by $\be_{i}$ and sum. The basis vectors are constants and pass
through both integrals; on the right, $u_{i}(\ip\bv\bn)\be_{i}
=(\ip\bv\bn)\bu=(\bu\otimes\bv)\bn$ by the definition of the tensor
product. Hence, with the caveat of Remark~\ref{rem:symbolic} still in
force,
\begin{equation}
   \frac{\dl}{\dl t}\int_{P_{t}}\bu\,\dl v
   =\int_{P_{t}}\bu'\,\dl v
   +\int_{\partial P_{t}}(\bu\otimes\bv)\bn\,\dl a .
   \tag{5.18}\label{eq:5.18}
\end{equation}

%-----------------------------------------------------------------------
\subsection{The flux derivative}\label{sec:fluxder}
%-----------------------------------------------------------------------

Balance laws for electromagnetic fields need the rate of change not of a
volume integral but of the \emph{flux} of a vector field through a moving
material surface. The claim is that
\begin{equation}
   \frac{\dl}{\dl t}\int_{\partial P_{t}}\ip\bu\bn\,\dl a
   =\int_{\partial P_{t}}\ip{\bu^{*}}\bn\,\dl a
   \tag{5.19}\label{eq:5.19}
\end{equation}
for a field $\bu^{*}$, the \emph{flux derivative} of $\bu$, to be
determined. As before the derivative follows a fixed set of material
points, here the ones occupying the evolving surface $\partial P_{t}$, and
as before we transform to a fixed domain before differentiating.

\begin{remark}[why the cofactor, and not $\bF$]\label{rem:whycof}
The pull-back of Remark~\ref{rem:pullback} needs a conversion factor for
whatever measure is being integrated, and the factor is different for each
of the three. Chapter 2 produced all three, and this chapter uses all
three:
\[
\begin{aligned}
   \text{line:}\quad &\dl\bx=\bF\,\dl\bX
     &&\text{(the definition of }\bF\text{),}\\
   \text{area:}\quad &\bn\,\dl a=\bF^{*}\bN\,\dl A
     &&\text{(Piola--Nanson, \eqr{2.54}),}\\
   \text{volume:}\quad &\dl v=J\,\dl V
     &&\text{(}J=\det\bF\text{).}
\end{aligned}
\]
$\bF$ maps a line element because a line element \emph{is} a vector
attached to a material fibre. An area element is not that kind of object:
$\bn\,\dl a$ is built from a \emph{pair} of fibres through their cross
product, and the tensor that carries cross products is by definition the
cofactor, $\bF^{*}(\ba\times\bb)=\bF\ba\times\bF\bb$
(Remark~\ref{rem:cofgeom}). Feeding $\bn$ to $\bF$ would deform the normal
as if it were a fibre, and normals are not fibres. Under a simple shear a normal tilts the
opposite way to a fibre lying in the shearing plane. That is the whole reason $\bF^{*}$ exists, and it is
why every surface integral in this chapter goes through Piola--Nanson
while every volume integral goes through $J$.

The identity $\bF^{*}=J\bF^{-t}$ of Proposition~\ref{prop:I2cof} then says
these are not three unrelated facts: the area factor is the volume factor
divided by the line factor, transposed, which is dimensionally exactly
what one would guess from ``area $=$ volume per length''.
\end{remark}

Piola--Nanson, $\bn\,\dl a=\bF^{*}\bN\,\dl A$, moves the integral to
$\partial P$; then $\bF^{*}=J\bF^{-t}$ (Proposition~\ref{prop:I2cof})
lets us shift $\bF$ across the inner product,
$\ip\bu{J\bF^{-t}\bN}=J\ip{\bF^{-1}\bu}\bN$:
\begin{equation}
\begin{aligned}
   \frac{\dl}{\dl t}\int_{\partial P_{t}}\ip\bu\bn\,\dl a
   &=\frac{\dl}{\dl t}\int_{\partial P}\ip\bu{\bF^{*}\bN}\,\dl A
    =\frac{\dl}{\dl t}\int_{\partial P}\ip{J\bF^{-1}\bu}\bN\,\dl A\\
   &=\int_{\partial P}\ip{\big(J\bF^{-1}\bu\big)^{\textstyle\cdot}}\bN\,\dl A\\
   &=\int_{\partial P}\ip{\big[J\bF^{-1}\dot\bu+\dot J\bF^{-1}\bu
     +J\big(\bF^{-1}\big)^{\textstyle\cdot}\bu\big]}\bN\,\dl A .
\end{aligned}
   \tag{5.20}\label{eq:5.20}
\end{equation}

The middle factor is reduced by a trick already used in Problem 1(d) of
Chapter 4. Differentiate $\bF\bF^{-1}=\I$:
\[
   \dot\bF\bF^{-1}+\bF\big(\bF^{-1}\big)^{\textstyle\cdot}=\Ob
   \quad\Longrightarrow\quad
   \bF\big(\bF^{-1}\big)^{\textstyle\cdot}=-\dot\bF\bF^{-1}=-\bL ,
\]
using $\bL=\dot\bF\bF^{-1}$ from \eqr{4.3}. To free
$\big(\bF^{-1}\big)^{\textstyle\cdot}$, multiply on the left by
$\bF^{-1}$:
\[
   \bF^{-1}\bF\big(\bF^{-1}\big)^{\textstyle\cdot}
   =\big(\bF^{-1}\big)^{\textstyle\cdot}
   =-\bF^{-1}\bL .
\]
Now substitute this and $\dot J=J\tr\bL$ into the bracket of \eqr{5.20},
taking the three terms in order:
\[
\begin{aligned}
   J\bF^{-1}\dot\bu
     &\quad\text{stays as it is,}\\
   \dot J\bF^{-1}\bu
     &=J(\tr\bL)\bF^{-1}\bu
      =J\bF^{-1}(\tr\bL)\bu ,\\
   J\big(\bF^{-1}\big)^{\textstyle\cdot}\bu
     &=J\big(-\bF^{-1}\bL\big)\bu
      =-J\bF^{-1}\bL\bu .
\end{aligned}
\]
In the second line the scalar $\tr\bL$ was moved across $\bF^{-1}$, which
is allowed because scalars commute with everything. Every term now carries
the same left factor $J\bF^{-1}$, so it can be taken out:
\[
   J\bF^{-1}\dot\bu+J\bF^{-1}(\tr\bL)\bu-J\bF^{-1}\bL\bu
   =J\bF^{-1}\big[\dot\bu+(\tr\bL)\bu-\bL\bu\big].
\]
That was the point of the whole manipulation. The bracket is a single
spatial vector, and the factor in front is exactly the one that converts
back to the current configuration. Reversing the step that produced it,
$\ip{J\bF^{-1}\bs}\bN=\ip\bs{J\bF^{-t}\bN}=\ip\bs{\bF^{*}\bN}$ for any
$\bs$, and then reading $\bF^{*}\bN\,\dl A$ as $\bn\,\dl a$ once more,
\begin{equation}
\begin{aligned}
   \frac{\dl}{\dl t}\int_{\partial P_{t}}\ip\bu\bn\,\dl a
   &=\int_{\partial P}\ip{\big[\dot\bu+(\tr\bL)\bu-\bL\bu\big]}{\bF^{*}\bN}\,\dl A
   \\
   &=\int_{\partial P_{t}}\ip{\big[\dot\bu+(\tr\bL)\bu-\bL\bu\big]}\bn\,\dl a .
\end{aligned}
   \tag{5.21}\label{eq:5.21}
\end{equation}
Comparison with \eqr{5.19} identifies $\bu^{*}=\dot\bu+(\dvg\bv)\bu
-\bL\bu$. Finally trade the material derivative for the spatial one with
\eqr{2.20}, $\dot\bu=\bu'+(\grad\bu)\bv$, and note that
$(\grad\bu)\bv+(\dvg\bv)\bu=\dvg(\bu\otimes\bv)$ by
Proposition~\ref{prop:divprod}:
\begin{equation}
   \boxed{\ \bu^{*}=\bu'+\dvg(\bu\otimes\bv)-(\grad\bv)\bu .\ }
   \tag{5.22}\label{eq:5.22}
\end{equation}

\begin{remark}[why $\bu^{*}$ is not $\dot\bu$]\label{rem:fluxwhy}
Three things change when a material surface is carried along by the
motion: the field on it, its area, and its orientation. The first term of
$\bu^{*}$ accounts for the field, the second for the stretching of the
area element, and the third, $-\bL\bu$, for the tilting of $\bn$. Drop the
motion, $\bv=\bze$, and $\bu^{*}$ collapses to $\bu'$ as it must. Problem 2
gives $\bu^{*}$ a second form in which the same three effects are
repackaged as a curl, which is how it is usually met in electrodynamics.
\end{remark}

%-----------------------------------------------------------------------
\subsection{Non-material regions}\label{sec:nonmat}
%-----------------------------------------------------------------------

Reynolds' formula assumed the region was material, meaning that it pulls
back to a fixed $P\subset\kappa$. Control volumes in fluid mechanics do not, and
neither will the shock surfaces of Section~\ref{sec:ch5disc}. Let
$V(t)\subset\kappa_{t}$ be a region whose boundary moves with local normal
speed $u(\bx,t)$ in the direction of its exterior unit normal $\bn$, with
$\varphi$ defined on $V$. The formula to be proved is
\begin{equation}
   \frac{\dl}{\dl t}\int_{V}\varphi\,\dl v
   =\int_{V}\varphi'\,\dl v+\int_{\partial V}\varphi\,u\,\dl a .
   \tag{5.23}\label{eq:5.23}
\end{equation}
It reduces to \eqr{5.15} when $V$ is material, for then the boundary is
carried by the motion and $u=\ip\bv\bn$.

The pull-back of Remark~\ref{rem:pullback} is unavailable here, because
there is no fixed domain to pull back to. We therefore go back to the
definition of the derivative:
\begin{equation}
   \frac{\dl}{\dl t}\int_{V(t)}\varphi\,\dl v
   =\lim_{\eps\to0}\frac1\eps\left[
     \int_{V(t+\eps)}\tilde\varphi(\bx,t+\eps)\,\dl v
     -\int_{V(t)}\tilde\varphi(\bx,t)\,\dl v\right].
   \tag{5.24}\label{eq:5.24}
\end{equation}
Two things differ between the two integrals in \eqr{5.24}: the integrand
has moved from time $t$ to time $t+\eps$, and so has the domain. A
difference in which two things changed at once is hard to estimate, so we
separate them by the standard device of adding and subtracting an
intermediate term. The intermediate we want is the one that agrees with
the first integral in its integrand and with the second in its domain,
namely
\[
   \int_{V(t)}\tilde\varphi(\bx,t+\eps)\,\dl v ,
\]
the integral over the \emph{old} domain of the \emph{new} integrand.
Adding and subtracting it inside the bracket of \eqr{5.24} costs nothing,
since the two copies cancel, and it splits the bracket into two
differences in each of which only one thing has changed. The first
difference has the same integrand on two different domains; the second has
two different integrands on the same domain. Writing the limit of the sum
as the sum of the limits, which is legitimate once both limits are known
to exist,
\begin{equation}
\begin{aligned}
   \frac{\dl}{\dl t}\int_{V(t)}\varphi\,\dl v
   &=\lim_{\eps\to0}\frac1\eps\left[
     \int_{V(t+\eps)}\tilde\varphi(\bx,t+\eps)\,\dl v
     -\int_{V(t)}\tilde\varphi(\bx,t+\eps)\,\dl v\right]\\
   &\quad+\lim_{\eps\to0}\frac1\eps\left[
     \int_{V(t)}\tilde\varphi(\bx,t+\eps)\,\dl v
     -\int_{V(t)}\tilde\varphi(\bx,t)\,\dl v\right]\\
   &=\lim_{\eps\to0}\frac1\eps
     \int_{V(t+\eps)\setminus V(t)}\tilde\varphi(\bx,t+\eps)\,\dl v
     \;+\;\int_{V(t)}\varphi'\,\dl v .
\end{aligned}
   \tag{5.25}\label{eq:5.25}
\end{equation}
Both simplifications in the last line deserve a word.

The second limit is now an ordinary derivative. The domain $V(t)$ is held
fixed while $\eps$ runs to zero, so the limit may be taken inside the
integral, and what is left inside is
\[
   \lim_{\eps\to0}\frac{\tilde\varphi(\bx,t+\eps)-\tilde\varphi(\bx,t)}{\eps}
   =\frac{\partial\tilde\varphi}{\partial t}\Big|_{\bx}=\varphi' ,
\]
the derivative at fixed place, by definition. This is where the prime of
Remark~\ref{rem:twoclocks} enters the formula, and it enters for a plain
reason: the domain was frozen at $V(t)$ for this piece, so the only thing
still moving was the field, and a field watched at a frozen place is
watched at fixed $\bx$.

The first limit simplifies because the two domains overlap. The region
$V(t)$ sits inside $V(t+\eps)$ when the boundary is advancing, so the
integral over $V(t)$ cancels the part of the integral over $V(t+\eps)$
taken over the same points, and what survives is the integral over the
thin shell $V(t+\eps)\setminus V(t)$ between the two boundaries. Where the
boundary is receding the shell counts with the opposite sign, which is
handled automatically once $u$ is allowed to be negative there. That
shell integral is the whole content of the formula, and evaluating it is
the business of the next paragraph.

\bookfig{1}\begin{figure}[htbp]\centering
\begin{tikzpicture}[scale=1.0,
   ar/.style={-{Stealth[length=2.2mm]},thick}]
% region at time t
\draw[thick,fill=blue!6]
  plot[smooth cycle,tension=0.85] coordinates
  {(0,1.15) (1.30,0.75) (1.55,-0.45) (0.55,-1.25) (-0.85,-1.00) (-1.35,0.10)};
% region at time t+eps: slightly inflated
\draw[thick,dashed,red!65!black]
  plot[smooth cycle,tension=0.85] coordinates
  {(0.06,1.55) (1.66,0.95) (1.92,-0.58) (0.66,-1.58) (-1.10,-1.28) (-1.70,0.12)};
\node[font=\small] at (-0.15,0.05) {$V(t)$};
\node[font=\small,text=red!65!black] at (1.62,1.35) {$V(t+\eps)$};
% the swept shell
\node[font=\footnotesize,align=center,text=gray!85] at (3.75,0.95)
  {shell $V(t+\eps)\setminus V(t)$\\ thickness $u\eps+o(\eps)$};
\draw[gray!70,-{Stealth[length=1.8mm]}] (3.05,0.72) .. controls (2.35,0.55) .. (1.72,0.28);
% normal and speed at a boundary point
\coordinate (P) at (1.52,-0.05);
\draw[ar] (P) -- ++(0.62,0.05) node[right,font=\small] {$\bn$};
\draw[ar,red!65!black] (P) -- ++(0.36,0.03)
   node[below right=-1pt,font=\small,text=red!65!black] {$u\eps$};
\fill (P) circle (1.1pt);
% area element
\draw[thick] (1.44,-0.42) -- (1.58,0.32);
\node[font=\small,anchor=north east] at (1.50,-0.40) {$\dl a$};
\end{tikzpicture}
\caption{Evolution of a non-material region. In time $\eps$ the boundary
advances a distance $u\eps$ along its own normal, so the shell swept out
has volume measure $\dl v=u\eps\,\dl a+o(\eps)$ over each element $\dl a$
of $\partial V(t)$. Dividing by $\eps$ and letting $\eps\to0$ leaves
$\int_{\partial V}\varphi u\,\dl a$: only the normal motion of the
boundary contributes, since sliding the boundary along itself sweeps no
volume.}
\label{fig:ch5nonmat}
\end{figure}

The shell is thin, and that is what makes it computable. Take one element
$\dl a$ of $\partial V(t)$ and ask how much volume the boundary sweeps
over it in time $\eps$. The boundary moves with normal speed $u$, so in
time $\eps$ it advances a distance $u\eps$ along $\bn$. The piece of shell
standing over $\dl a$ is therefore a slab of base $\dl a$ and height
$u\eps$, of volume $u\eps\,\dl a$, together with a correction.

The correction is $o(\eps)$, meaning it vanishes faster than $\eps$ does,
and it has two sources. The boundary need not move exactly along $\bn$;
but by the argument after \eqr{5.27} the tangential part sweeps no volume
to first order, since sliding a surface along itself moves no material
across it. And the surface curves, so the slab is not exactly a cylinder;
but the deviation is of order $\eps^{2}$, being a product of the advance
$u\eps$ with the change in $\bn$ over that advance, itself of order
$\eps$. Both corrections therefore disappear on dividing by $\eps$ and
letting $\eps$ tend to zero, which is the only thing we will do with them.

So, with reference to Figure~\ref{fig:ch5nonmat}, the volume measure on
the shell is
\[
   \dl v=u\eps\,\dl a+o(\eps),
\]
and the integral over the shell becomes an integral over $\partial V(t)$.
Therefore
\begin{equation}
\begin{aligned}
   \lim_{\eps\to0}\frac1\eps
   \int_{V(t+\eps)\setminus V(t)}\tilde\varphi(\bx,t+\eps)\,\dl v
   &=\lim_{\eps\to0}\int_{\partial V(t)}
     \tilde\varphi(\bx,t+\eps)\,u\,\dl a+\frac{o(\eps)}{\eps}
   \\
   &=\int_{\partial V(t)}\tilde\varphi(\bx,t)\,u\,\dl a ,
\end{aligned}
   \tag{5.26}\label{eq:5.26}
\end{equation}
which with \eqr{5.25} gives \eqr{5.23}.

It is convenient to carry a vector field $\bu(\bx,t)$ on $\partial V(t)$
with $\ip\bu\bn=u$. Only the normal component has any meaning: the
tangential part $(\I-\bn\otimes\bn)\bu$ slides the boundary along itself
and sweeps no volume, so it cannot appear in \eqr{5.23}. Then
\begin{equation}
   \frac{\dl}{\dl t}\int_{V}\varphi\,\dl v
   =\int_{V}\varphi'\,\dl v+\int_{\partial V}\varphi\,\ip\bu\bn\,\dl a ,
   \tag{5.27}\label{eq:5.27}
\end{equation}
in exact analogy with \eqr{5.15}. The analogy is only formal, however:
$\bv$ is defined throughout the body, $\bu$ only on $\partial V$.

%=======================================================================
\section{Discontinuous fields}\label{sec:ch5disc}
%=======================================================================

%-----------------------------------------------------------------------
\subsection{Reynolds' formula for fields that suffer a discontinuity
across a surface}\label{sec:reydisc}
%-----------------------------------------------------------------------

Shock waves, phase boundaries and propagating cracks all present the same
picture: a surface $s(t)$ sweeping through a material volume
$P_{t}\subset\kappa_{t}$ at speed $u$ along its unit normal $\bn$, cutting
$P_{t}$ into two parts $V^{\pm}$ across which the fields jump. The two
parts are \emph{not} material regions, and $s$ is not a material surface,
whenever $u\neq\ip\bv\bn$ (Figure~\ref{fig:ch5shock}).

Let $\varphi$ be discontinuous across $s$, which we then call a
\emph{singular surface}. Its jump at $\bx\in s$ is
\begin{equation}
   \jump{\varphi}=\varphi^{+}-\varphi^{-},
   \tag{5.28}\label{eq:5.28}
\end{equation}
where $\varphi^{\pm}$ are the limits of $\varphi$ as $\bx$ is approached
from the interiors of $V^{\pm}$. Throughout, $\bn$ points from $V^{-}$
into $V^{+}$; the sign of every jump depends on that convention and on
nothing else. From the figure,
\begin{equation}
   P_{t}=V^{+}\cup V^{-},\qquad
   \partial P_{t}^{\pm}=\partial V^{\pm}\cap\partial P_{t},\qquad
   \partial P_{t}=\partial P_{t}^{+}\cup\partial P_{t}^{-}.
   \tag{5.29}\label{eq:5.29}
\end{equation}

\bookfig{2}\begin{figure}[htbp]\centering
\begin{tikzpicture}[scale=1.05,
   ar/.style={-{Stealth[length=2.2mm]},thick}]
% material region
\draw[thick] plot[smooth cycle,tension=0.82] coordinates
  {(0,1.55) (1.75,1.05) (2.05,-0.55) (0.70,-1.70) (-1.15,-1.35) (-1.80,0.15)};
% the propagating surface s : a curve cutting it
\draw[very thick,red!65!black] (-1.62,-0.72) .. controls (-0.25,-0.30) and (0.55,0.30) .. (1.92,0.28);
\node[font=\small,text=red!65!black] at (2.28,0.52) {$s(t)$};
% shading labels
\node[font=\small] at (-0.35,0.95) {$V^{+}$};
\node[font=\small] at (0.15,-1.05) {$V^{-}$};
\node[font=\footnotesize,text=gray!80,anchor=west] at (2.20,1.30)
  {$\partial P^{+}_{t}$};
\draw[gray!70,-{Stealth[length=1.6mm]}] (2.18,1.24) -- (1.35,1.16);
\node[font=\footnotesize,text=gray!80,anchor=west] at (2.20,-1.35)
  {$\partial P^{-}_{t}$};
\draw[gray!70,-{Stealth[length=1.6mm]}] (2.18,-1.29) -- (1.25,-1.05);
% normal on s, pointing from V- into V+
\coordinate (Q) at (0.20,0.06);
\draw[ar] (Q) -- ++(-0.16,0.72) node[above,font=\small] {$\bn$};
\draw[ar,red!65!black] (Q) -- ++(-0.10,0.44)
   node[right=1pt,font=\small,text=red!65!black] {$u$};
\fill (Q) circle (1.2pt);
% material velocity, deliberately not aligned with u n
\draw[ar,blue!55!black] (Q) -- ++(0.62,0.30)
   node[right,font=\small,text=blue!55!black] {$\bv$};
\end{tikzpicture}
\caption{A material region $P_{t}\subset\kappa_{t}$ traversed by a
propagating surface $s$. The normal $\bn$ on $s$ points from $V^{-}$ into
$V^{+}$ and $s$ advances at speed $u$ along it. Material crosses $s$
whenever $u\neq\ip\bv\bn$: the surface is then not made of the same
particles from one instant to the next, and $V^{\pm}$ are non-material
regions even though their union $P_{t}$ is material. This is why
\eqr{5.27}, and not \eqr{5.15}, is the formula to apply to each half.}
\label{fig:ch5shock}
\end{figure}

Let $\varphi$ be smooth in $V^{+}$ and in $V^{-}$ separately. Each half is
a non-material region, so \eqr{5.27} applies to each. The boundary of
$V^{+}$ consists of $\partial P_{t}^{+}$, which is material and therefore
moves with normal speed $\ip\bv\bn$, and of $s$, whose exterior normal
\emph{as seen from $V^{+}$} is $-\bn$, so that its normal speed in that
direction is $-u$. For $V^{-}$ the exterior normal on $s$ is $+\bn$ and
the speed is $+u$. Hence
\begin{equation}
\begin{aligned}
   \frac{\dl}{\dl t}\int_{V^{+}}\varphi\,\dl v
   &=\int_{V^{+}}\varphi'\,\dl v
    +\int_{\partial P^{+}_{t}}\varphi\,\ip\bv\bn\,\dl a
    +\int_{s}\varphi^{+}(-u)\,\dl a
   \qquad\text{and}\\[1mm]
   \frac{\dl}{\dl t}\int_{V^{-}}\varphi\,\dl v
   &=\int_{V^{-}}\varphi'\,\dl v
    +\int_{\partial P^{-}_{t}}\varphi\,\ip\bv\bn\,\dl a
    +\int_{s}\varphi^{-}u\,\dl a .
\end{aligned}
   \tag{5.30}\label{eq:5.30}
\end{equation}
Before adding them, notice why the signs came out as they did, since this
is the only delicate point in the section. Equation \eqr{5.27} was derived
for a region whose boundary advances at speed $u$ \emph{along its own
exterior normal}. The surface $s$ is part of the boundary of both halves,
but the two halves look at it from opposite sides. Seen from $V^{-}$, the
outward direction is $+\bn$ and the surface advances into $V^{+}$, which
is away from $V^{-}$: the region $V^{-}$ is gaining, and the speed to use
is $+u$. Seen from $V^{+}$, the outward direction is $-\bn$ and the
surface is advancing into $V^{+}$ itself: the region is losing, and the
speed to use is $-u$. One surface, one physical motion, two opposite
bookkeeping signs, because the two regions are on opposite sides of it.
The limiting value of $\varphi$ likewise differs, being $\varphi^{+}$ from
one side and $\varphi^{-}$ from the other, which is exactly what makes the
sum fail to cancel.

Now add the two lines of \eqr{5.30}, one piece at a time. The volume
integrals combine, because $V^{+}\cup V^{-}=P_{t}$ by \eqr{5.29}$_{1}$ and
the two parts overlap only on $s$, a surface of zero volume:
\[
   \int_{V^{+}}\varphi'\,\dl v+\int_{V^{-}}\varphi'\,\dl v
   =\int_{P_{t}}\varphi'\,\dl v .
\]
The material surface integrals combine for the same reason, using
$\partial P_{t}=\partial P^{+}_{t}\cup\partial P^{-}_{t}$ from
\eqr{5.29}$_{3}$:
\[
   \int_{\partial P^{+}_{t}}\varphi\,\ip\bv\bn\,\dl a
   +\int_{\partial P^{-}_{t}}\varphi\,\ip\bv\bn\,\dl a
   =\int_{\partial P_{t}}\varphi\,\ip\bv\bn\,\dl a .
\]
The two integrals over $s$ are both over the same surface, so they combine
into one integrand:
\[
   \int_{s}\varphi^{+}(-u)\,\dl a+\int_{s}\varphi^{-}u\,\dl a
   =-\int_{s}\big(\varphi^{+}-\varphi^{-}\big)u\,\dl a
   =-\int_{s}\jump{\varphi}\,u\,\dl a ,
\]
the last step being the definition \eqr{5.28} of the jump. Collecting the
three results,
\begin{equation}
   \boxed{\ \frac{\dl}{\dl t}\int_{P_{t}}\varphi\,\dl v
   =\int_{P_{t}}\varphi'\,\dl v
   +\int_{\partial P_{t}}\varphi\,\ip\bv\bn\,\dl a
   -\int_{s}\jump{\varphi}u\,\dl a .\ }
   \tag{5.31}\label{eq:5.31}
\end{equation}

\begin{remark}[the sign, and what the extra term is]\label{rem:jumpsign}
The minus sign is not a convention one may take or leave: it follows from
$\bn$ pointing into $V^{+}$, which is also the convention that defines
$\jump\varphi$ in \eqr{5.28}. Reverse $\bn$ and both $\jump\varphi$ and
$u$ change sign, so the product, and the formula, are unchanged, as they must be, since the left-hand side knows nothing
about our choice. As for the
meaning: if $\varphi$ is smooth then $\jump\varphi=0$ and \eqr{5.31} is
\eqr{5.15}. When it is not, the surface acts as a source. It is this term
that will produce the Rankine--Hugoniot jump conditions in Chapter 6, in
exactly the way \eqr{5.15} produces the smooth balance laws.
\end{remark}

The passage to vector fields is the one already made from \eqr{5.15} to
\eqr{5.18}: apply \eqr{5.31} to each Cartesian component and reassemble.

%-----------------------------------------------------------------------
\subsection{The divergence theorem for discontinuous fields}
\label{sec:divdisc}
%-----------------------------------------------------------------------

The divergence theorem \eqr{5.1} needs the same repair. Recall that in
$\ip\ba\bn$ on its right-hand side, $\ba$ means the limit taken as
$\partial R$ is approached \emph{from inside} $R$. Let $\ba(\bx,t)$ be
smooth in $V^{+}$ and $V^{-}$ separately and suffer
\begin{equation}
   \jump\ba=\ba^{+}-\ba^{-}\qquad\text{at }\bx\in s .
   \tag{5.32}\label{eq:5.32}
\end{equation}

Apply \eqr{5.1} to each half separately. This is legitimate because
$\ba$ is smooth on each side taken by itself. The boundary of $V^{-}$ is $\partial P^{-}_{t}$ together
with $s$, on which the exterior normal is $+\bn$ and the relevant limit is
$\ba^{-}$; the boundary of $V^{+}$ is $\partial P^{+}_{t}$ together with
$s$, on which the exterior normal is $-\bn$ and the limit is $\ba^{+}$:
\begin{equation}
\begin{aligned}
   \int_{\partial P^{-}_{t}}\ip\ba\bn\,\dl a
   +\int_{s}\ip{\ba^{-}}\bn\,\dl a
   &=\int_{V^{-}}\dvg\ba\,\dl v
   \qquad\text{and}\\[1mm]
   \int_{\partial P^{+}_{t}}\ip\ba\bn\,\dl a
   +\int_{s}\ip{\ba^{+}}{(-\bn)}\,\dl a
   &=\int_{V^{+}}\dvg\ba\,\dl v .
\end{aligned}
   \tag{5.33}\label{eq:5.33}
\end{equation}
The sign pattern is the same one that appeared in \eqr{5.30} and for the
same reason: the two halves see the shared surface from opposite sides, so
one of them reads its exterior normal as $+\bn$ and the other as $-\bn$.
There is no motion in this argument at all, only geometry, which is why
$u$ never appears.

Add the two lines. The right sides combine into a single integral over
$P_{t}$, since $V^{+}\cup V^{-}=P_{t}$:
\[
   \int_{V^{-}}\dvg\ba\,\dl v+\int_{V^{+}}\dvg\ba\,\dl v
   =\int_{P_{t}}\dvg\ba\,\dl v .
\]
On the left, the two integrals over the outer boundary combine into one
over $\partial P_{t}$ by \eqr{5.29}$_{3}$, and the two over $s$ give
\[
   \int_{s}\ip{\ba^{-}}\bn\,\dl a+\int_{s}\ip{\ba^{+}}{(-\bn)}\,\dl a
   =-\int_{s}\ip{\big(\ba^{+}-\ba^{-}\big)}\bn\,\dl a
   =-\int_{s}\ip{\jump\ba}\bn\,\dl a ,
\]
using the linearity of the inner product in its first slot and then the
definition \eqr{5.32}. So the sum reads
\[
   \int_{\partial P_{t}}\ip\ba\bn\,\dl a-\int_{s}\ip{\jump\ba}\bn\,\dl a
   =\int_{P_{t}}\dvg\ba\,\dl v ,
\]
and moving the jump term to the other side,
\begin{equation}
   \boxed{\ \int_{\partial P_{t}}\ip\ba\bn\,\dl a
   =\int_{P_{t}}\dvg\ba\,\dl v
   +\int_{s}\ip{\jump\ba}\bn\,\dl a .\ }
   \tag{5.34}\label{eq:5.34}
\end{equation}

In every formula of this section the jump is taken with respect to a
normal $\bn$ on $s$ directed from $V^{-}$ to $V^{+}$.

\begin{remark}[a way to remember both]\label{rem:distrib}
Read \eqr{5.34} backwards: $\int_{P_{t}}\dvg\ba\,\dl v$ \emph{under-counts}
the boundary flux by exactly the flux the jump carries across $s$.
Equivalently, the distributional divergence of a piecewise smooth field
is the classical divergence plus a surface layer of strength
$\ip{\jump\ba}\bn$ on $s$. The same reading applies to \eqr{5.31}: the
singular surface carries its own share of the balance, and localizing an
integral statement over a region cut by $s$ therefore produces two local statements: a field equation away from $s$,
and a jump condition on $s$ itself.
\end{remark}

%=======================================================================
\setcounter{problemenv}{0}
\section{Problems}\label{sec:problems:c5}
%=======================================================================

%-----------------------------------------------------------------------
\begin{problem}[Transport of a line integral]
Fill in the missing terms in the transport formula
\[
   \frac{\dl}{\dl t}\int_{C_{t}}\ip\ba{\dl\bx}
   =\int_{C_{t}}\ip{\big(\ba'+\dots\big)}{\dl\bx},
\]
where $C_{t}$ is a moving material curve and $\ba'$ is the time derivative
of the smooth spatial field $\ba$ at fixed $\bx$.
\end{problem}

\begin{solution}
The pattern of the chapter applies unchanged: convert to a fixed domain,
differentiate, convert back. The fixed domain here is the material curve
$C\subset\kappa$ of which $C_{t}=\bchi_{\kappa}(C,t)$ is the image.

A material line element obeys $\dl\bx=\bF\,\dl\bX$ by the definition of
the deformation gradient, and $\dl\bX$ is fixed in time, so by \eqr{4.3}
\[
   \big(\dl\bx\big)^{\textstyle\cdot}=\dot\bF\,\dl\bX=\bL\bF\,\dl\bX
   =\bL\,\dl\bx .
\]
This is the line-element counterpart of $\dot J=J\tr\bL$ for the volume
element and of Piola--Nanson for the area element; the three together are
the transport rules for the three measures.

Now differentiate the integral over the fixed $C$ and use the product
rule:
\[
\begin{aligned}
   \frac{\dl}{\dl t}\int_{C_{t}}\ip\ba{\dl\bx}
   &=\int_{C}\Big[\ip{\dot\ba}{\dl\bx}
     +\ip\ba{\big(\dl\bx\big)^{\textstyle\cdot}}\Big]
    =\int_{C_{t}}\Big[\ip{\dot\ba}{\dl\bx}+\ip\ba{\bL\,\dl\bx}\Big]\\
   &=\int_{C_{t}}\ip{\big(\dot\ba+\bL^{t}\ba\big)}{\dl\bx},
\end{aligned}
\]
The last step moved $\bL$ from the second slot of the inner product to the
first. That is the definition of the transpose,
$\ip\ba{\bL\bs}=\ip{\bL^{t}\ba}\bs$, applied with $\bs=\dl\bx$, and then
the two terms were collected under one inner product by linearity in the
first slot.

Now trade the material derivative for the spatial one. By \eqr{2.20},
$\dot\ba=\ba'+(\grad\ba)\bv$, and $\bL=\grad\bv$ so $\bL^{t}
=(\grad\bv)^{t}$. Substituting both,
\[
   \boxed{\ \frac{\dl}{\dl t}\int_{C_{t}}\ip\ba{\dl\bx}
   =\int_{C_{t}}\ip{\Big(\ba'+(\grad\ba)\bv+(\grad\bv)^{t}\ba\Big)}{\dl\bx}.\ }
\]
So the missing terms are $(\grad\ba)\bv+(\grad\bv)^{t}\ba$. The first is
the convective change of $\ba$, the change met by being carried to a new
place. The second has a different origin: it is the change caused by the
line element itself rotating and stretching, and it would be absent if one
guessed the answer to be $\int\ip{\dot\ba}{\dl\bx}$. Guessing that is the
natural mistake, and the reason it fails is that $\dl\bx$ is not a passive
label but a material object that the motion deforms, exactly as $\dl v$
and $\dl a$ are.

\emph{A check worth doing.} Take $\ba=\bv$, so that the integral is the
circulation of the velocity around the curve. The integrand becomes
\[
   \ip{\bv'}{\dl\bx}+\ip{(\grad\bv)\bv}{\dl\bx}
   +\ip{(\grad\bv)^{t}\bv}{\dl\bx}.
\]
The third term simplifies. In components,
$\big[(\grad\bv)^{t}\bv\big]_{i}=v_{j,i}v_{j}
=\tfrac12\big(v_{j}v_{j}\big)_{,i}=\tfrac12\big(|\bv|^{2}\big)_{,i}$,
using the product rule backwards, so
\[
   (\grad\bv)^{t}\bv=\tfrac12\grad\big(|\bv|^{2}\big),
\]
a gradient. The line integral of a gradient around a closed curve is the
change of the potential around a closed loop, which is zero. So on a
closed material curve the third term drops out entirely and
\[
   \frac{\dl}{\dl t}\oint_{C_{t}}\ip\bv{\dl\bx}
   =\oint_{C_{t}}\ip{\big(\bv'+(\grad\bv)\bv\big)}{\dl\bx}
   =\oint_{C_{t}}\ip{\dot\bv}{\dl\bx},
\]
the last step recognising \eqr{2.21}. This is Kelvin's circulation theorem
in its purely kinematical form: the circulation around a closed material
curve changes only through the acceleration. Once Chapter 6 supplies an
equation of motion for $\dot\bv$, this becomes the statement that
circulation is conserved whenever the acceleration is itself a gradient,
which is the usual form of the theorem in fluid mechanics.
\end{solution}

%-----------------------------------------------------------------------
\begin{problem}[A second form of the flux derivative]
Show that the flux derivative of a vector field $\bu$ may be written
\[
   \bu^{*}=\bu'+(\dvg\bu)\bv-\crl(\bv\times\bu).
\]
\end{problem}

\begin{solution}
Start from the identity for the curl of a cross product, Problem 24 of
Chapter 1:
\[
   \crl(\ba\times\bb)=(\dvg\bb)\ba-(\dvg\ba)\bb
   +(\grad\ba)\bb-(\grad\bb)\ba .
\]
Put $\ba=\bv$, $\bb=\bu$:
\[
   \crl(\bv\times\bu)=(\dvg\bu)\bv-(\dvg\bv)\bu
   +(\grad\bv)\bu-(\grad\bu)\bv .
\]
Therefore
\[
\begin{aligned}
   (\dvg\bu)\bv-\crl(\bv\times\bu)
   &=(\dvg\bu)\bv-(\dvg\bu)\bv+(\dvg\bv)\bu\\
   &\qquad-(\grad\bv)\bu+(\grad\bu)\bv\\
   &=\underbrace{(\grad\bu)\bv+(\dvg\bv)\bu}
     _{=\ \dvg(\bu\otimes\bv)}-(\grad\bv)\bu ,
\end{aligned}
\]
the brace by Proposition~\ref{prop:divprod}. Adding $\bu'$ to both sides
and comparing with \eqr{5.22},
\[
   \boxed{\ \bu^{*}=\bu'+(\dvg\bu)\bv-\crl(\bv\times\bu).\ }
\]

The second form is the one that makes Faraday's law transparent. If
$\dvg\bu=0$, as it is for the magnetic induction, the middle term drops
and
\[
   \frac{\dl}{\dl t}\int_{\partial P_{t}}\ip\bu\bn\,\dl a
   =\int_{\partial P_{t}}\ip{\big[\bu'-\crl(\bv\times\bu)\big]}\bn\,\dl a :
\]
the flux through a material loop changes through the local rate $\bu'$
and through the motional term $-\crl(\bv\times\bu)$, whose line-integral
form $\oint\ip{(\bv\times\bu)}{\dl\bx}$ is the motional electromotive
force. The two forms \eqr{5.22} and the boxed one are of course the same
field; which is preferable depends on whether $\dvg\bu$ or
$\dvg(\bu\otimes\bv)$ is the quantity at hand.
\end{solution}

%-----------------------------------------------------------------------
\begin{problem}[The referential statements]
Everything in this chapter is written in the spatial description. Derive
the corresponding referential statements for fields suffering a
discontinuity across a surface $S$ propagating through a fixed region
$P\subset\kappa$ at speed $U$ along its unit normal $\bN$. \emph{Hint:}
use $\dot\bX=\bze$.
\end{problem}

\begin{solution}
The hint is the whole solution, once its consequence is seen. In the
referential description the independent variable is $\bX$, and material
points do not move in $\kappa$: $\dot\bX=\bze$. So the ``material
velocity'' of the referential picture vanishes, and every term in the
spatial formulas built from $\bv$ disappears.

\textbf{Transport.} Let $\hat\Phi(\bX,t)$ be a referential density and let
$S(t)$ propagate through the \emph{fixed} region $P$ at normal speed $U$
along $\bN$, cutting $P$ into $P^{\pm}$ with $\bN$ directed from $P^{-}$
into $P^{+}$. Apply the non-material formula \eqr{5.27} in $\kappa$ to
each half. Because $\partial P$ is fixed, its normal speed is zero and the
integral over $\partial P$ that replaced $\int_{\partial P_{t}}\varphi
\ip\bv\bn\,\dl a$ is absent; only the two integrals over $S$ survive, with
exterior normal $-\bN$ and speed $-U$ for $P^{+}$ and $+\bN$, $+U$ for
$P^{-}$. Adding,
\[
   \boxed{\ \frac{\dl}{\dl t}\int_{P}\hat\Phi\,\dl V
   =\int_{P}\frac{\partial\hat\Phi}{\partial t}\,\dl V
   -\int_{S}\jump{\hat\Phi}\,U\,\dl A .\ }
\]
Compare \eqr{5.31}: the boundary flux term is gone and the singular
surface term is untouched. That is the referential picture in one
line. Nothing is convected, because nothing moves; all that happens is that the
discontinuity sweeps across the material.

\textbf{Divergence.} The argument of Section~\ref{sec:divdisc} used only
the smoothness of the field on each side and the direction of the
exterior normal, never the motion. It therefore transcribes verbatim, with
$\Dvg$ for $\dvg$ and $\bN$, $\dl A$, $\dl V$ for $\bn$, $\dl a$, $\dl v$:
for a referential field $\bA(\bX,t)$ smooth on each side of $S$,
\[
   \boxed{\ \int_{\partial P}\ip\bA\bN\,\dl A
   =\int_{P}\Dvg\bA\,\dl V+\int_{S}\ip{\jump\bA}\bN\,\dl A .\ }
\]

\textbf{Relation between $U$ and $u$.} The two speeds describe the same
physical surface seen in two configurations and are not equal. A point of
$S$ with referential normal speed $U$ maps to a point of $s$; the
displacement $U\eps\bN$ in $\kappa$ carries over to $\bF(U\eps\bN)$ in
$\kappa_{t}$, whose component along $\bn$ is the normal advance there.
Using Piola--Nanson to relate the normals, one finds
\[
   u-\ip\bv\bn=\frac{J}{|\bF^{*}\bN|}\,U ,
\]
so the two vanish together: $S$ is material in $\kappa$ exactly when $s$
is material in $\kappa_{t}$, as it must be. The combination
$u-\ip\bv\bn$, the speed of the surface \emph{relative to the material},
is the one with physical content, and it is what Chapter 6 will meet in
the mass jump condition.
\end{solution}
